
\chapter{Orthogonality and Inner Products}

\subsection{Orthogonality and Projection}

Let \( a, b \in V \). The vectors \(a\) and \(b\) are orthogonal to each other if

\[
    \inner{a}{b} = 0.
\]

This is written as \( a \perp b \).

Also, the Pythagorean Theorem holds

\[
    \metric{a + b}^2 = \metric{a}^2 + \metric{b}^2,
\]

for \( a \perp b \) in all unitary vector spaces.

For the orthogonal projection \( p_b(a) \) of a vector \(a\) onto \(b\), with \( b \neq 0 \), the 
formula in every unitary vector space is:

\[
    p_b(a) = \frac{\inner{a}{b}}{\inner{b}{b}} b.
\]

\subsection{Induced Metric}

Given \((V, \metric{\cdot})\) a metric space there is an inner product \(\inner{\cdot}{\cdot}\) on \(\field\) such that 
\(\metric{x}^2 = \inner{x}{x}\) if and only if the \emph{parallelepiped equation} holds. 

\[
    \metric{x + y}^2 + \metric{x - y}^2 = 2(\metric{x}^2 + \metric{y}), \ \forall x,y \in V
\]

In such a case, the inner product is given by 

\[
    \inner{x}{y} = \frac{1}{4} (\metric{x + y}^2 - \metric{x - y}^2)
\]

\textbf{Proof:}

\(\Rightarrow\): If \(\|x\|^2 = \inner{x}{x}\) then 

\begin{align*} 
    \metric{x + y}^2 + \metric{x - y}^2 &= \inner{x + y}{x + y} + \inner{x - y}{x - y} \\ 
    &= 2 (\inner{x}{x} + \inner{y}{y}) \\ 
    \inner{x}{x + y} + \inner{y}{x + y} + \inner{x}{x - y} - \inner{y}{x - y} &= 2 (\inner{x}{x} + \inner{y}{y}) \\ 
    \inner{x}{x} + \inner{x}{y} + \inner{y}{y} + \inner{x}{y} + \inner{x}{x} - \inner{x}{y} - \inner{y}{x} + \inner{y}{y} &= 2 (\inner{x}{x} + \inner{y}{y}) \\ 
    2\inner{x}{x} +  2\inner{y}{y} +  &= 2 (\inner{x}{x} + \inner{y}{y})
\end{align*}

\(\Leftarrow\): The parallelepiped equation holds, then want to show that the equation 

\[
    f(x, y) = \frac{1}{4} (\metric{x + y}^2 - \metric{x - y}^2)
\]

Satisfies the properties of a inner product:

\begin{itemize}

    \item \(f(x,y) \ge 0\) trivial. 

    \item \(f(x,y) = 0 \iff x = y = 0\) is also trivial and

    \item \(f(x, y) = f(y,x)\) is also trivial by our definition.

    \item \(f(\alpha x + y, z) = \alpha f(x, y) + f(y, z)\) Must be shown in a two step process.

\end{itemize}

First we show that \(f(x + y, z) = f(x,z) + f(y,z)\):

\begin{align*}
    4(f(x,y) + f(z, y)) &= \metric{x + y}^2 - \metric{x - y}^2 + \metric{z + y}^2 - \metric{z - y}^2 \\ 
    &= \metric{x + y}^2 +\metric{z}^2 \\ 
    &- \metric{x - y}^2 - \metric{z}^2  \\ 
    &+ \metric{z + y}^2 + \metric{x}^2 \\ 
    &- \metric{z - y}^2 - \metric{x}^2 
\end{align*}

Hence, by adding this clever zero and using the parallelepiped equation we get

\begin{align*}
    2(f(x, y) + f(z, y)) &= \metric{x + y + z}^2 + \metric{x + y - z}^2 \\
    &- \metric{x - y + z}^2 - \metric{x - y - z}^2 \\
    &+ \metric{z + y + x}^2 + \metric{z + y - x}^2 \\
    & - \metric{z - y + x}^2 - \metric{z - y - x}^2 \\
    &= 2(\metric{x + z + y}^2 - \metric{x + z - y}^2) \\
    &= 2f(x + z, y)
\end{align*}

Therefore, \(f(x + z, y) = f(x,y) + f(z, y)\).

Now we need to prove that \(f(\alpha x, y) = \alpha f(x, y)\):

For the Natural starting with \(n = 2\): \(f(2x, y) = f(x + x, y) = f(x, y) + f(x, y) = 2f(x, y)\), hence 
using induction for \(n \in \N\) with base case 2 and summation, the prove is concluded.

For the Rationals, \(r \in \Q\) with \(r = \frac{m}{n}\) we have that 

\[
    f(rx, y) = f\left(\frac{mx}{n}, y\right) = m f\left(\frac{1}{n}x, y\right) =
    \frac{m}{n}n f\left(\frac{1}{n}x, y\right) = \frac{m}{n} f\left(x, y\right)
\]

Finally, for the Reals; each Real number can be approximated as a limit of an infinite sequence of Rationals.

For \(\alpha \in \R\) there is \(r_K \in \Q\) such that \(\lim_{k \to \infty} r_k = \alpha\). 

Given that \(f\) is continuous we can then use the following 

\[
    f\left(\lim_{k \to \infty} r_K x, y\right) = \lim_{k \to \infty}f(r_K x, y) = \lim_{k \to \infty}r_k f(x, y) = \alpha f(x, y)
\]

\QED

\subsection{Euclidean Inner Product as Sum of Metrics}

Given our Euclidean inner product for any two vector in \(\R^n\), the following holds:

\[
    \inner{x_i}{x_j} = \frac{1}{2} (\metric{x_i}^2 + \metric{x_j}^2 - \metric{x_i - x_j}^2)
\]

\textbf{Proof:}

\begin{align*} 
    \metric{x_i}^2 + \metric{x_j}^2 - \metric{x_i - x_j}^2 &= \inner{x_i}{x_i}+ \inner{x_j}{x_j} - \inner{x_i - x_j}{x_i - x_j} \\ 
    &= \sum (x_i^{(k)})^2 + \sum (x_i^{(k)})^2  - \sum (x_i^{(k)} - x_j^{(k)})^2 \\
    &= \sum (x_i^{(k)})^2 + \sum (x_i^{(k)})^2  - \left[\sum (x_i^{(k)})^2 - 2 \sum x_i^{(k)} x_j^{(k)} + \sum (x_j^{(k)})^2\right] \\ 
    &= \sum (x_i^{(k)})^2 + \sum (x_i^{(k)})^2  -\sum (x_i^{(k)})^2 + 2 \sum x_i^{(k)} x_j^{(k)} - \sum (x_j^{(k)})^2 \\ 
    &= 2 \sum x_i^{(k)} x_j^{(k)}
\end{align*}

Now by substituting into the original equality

\begin{align*} 
    \inner{x_i}{x_j} &= \frac{1}{2} (\metric{x_i}^2 + \metric{x_j}^2 - \metric{x_i - x_j}^2) \\ 
     &= \frac{1}{2} \left(2 \sum x_i^{(k)} x_j^{(k)}\right) \\ 
     &= \sum x_i^{(k)} x_j^{(k)} =  \inner{x_i}{x_j}
\end{align*}

\QED

\subsection{Orthogonal Projection and Orthogonal Complement}

Let \( U \) be a finitely generated subspace of \( V \) and \( a \in V \). A vector 
\( p_U(a) \in U \) is called the orthogonal projection of \(a\) onto \( U \) if

\[
    a - p_U(a) \perp u \quad \forall u \in U
\]

The question arises about well-definedness, i.e., whether such a vector \( p_U(a) \) always exists and 
if it is unique. The following concept helps in the discussion of uniqueness.

For \( M \subseteq V \), the \emph{orthogonal complement} of \( M \) is defined as

\[
    M^\perp = \{ v \in V \mid v \perp u \, \forall u \in M \}.
\]

\begin{itemize}
    \item \( M^\perp \) is a subspace of \( V \).
    \item Let \( U \) be a subspace of \( V \). Then, we have \( U \cap U^\perp = \{0\} \).
\end{itemize}

\textbf{Proof:}

We need to check the closure property. For \( u \in M \), \( x, y \in M^\perp \), and 
\( \lambda \in \R \), we have:

\[
    \inner{x + y}{u} = \inner{x}{u} + \inner{y}{u} = 0
\]
\[
    \inner{\lambda x}{u} = \lambda \inner{x}{u} = 0.
\]

Let \( a \in U \cap U^\perp \). Then, we have \( \inner{a}{u} = 0 \) for all \( u \in U \), 
because \( a \in U^\perp \), and in particular, \( \inner{a}{a} = 0 \) since \( a \in U \), 
which implies \( a = 0 \).

\subsection{Orthogonality of the basis to the Complement}

Let \( U \) be as before, and let \( (u_1, \ldots, u_m) \) be a basis of \( U \). For \( v \in V \), we 
have:

\[
    v \in U^\perp \quad \text{if and only if} \quad \inner{v}{u_i} = 0 \quad \forall 1 \leq i \leq m.
\]

\subsection{How to find the orthogonal complement}

Let \( U \) be a finitely generated subspace of \( V \) with basis \( (u_1, \ldots, u_m) \). To find the 
orthogonal complement \( U^\perp \), we can solve the system of equations:

\[
\inner{v}{u_i} = 0 \quad \forall 1 \leq i \leq m.
\]

This system can be expressed in matrix form as \( A \cdot v = 0 \), where \(A\) is the matrix whose rows 
are the vectors \( u_i \) and \( v \) is the vector we want to find in \( U^\perp \).

\subsection{Orthogonal and Orthonormal Systems}

Let \( B = (v_1, \ldots, v_m) \) be an \( m \)-tuple of vectors in \( V \setminus \{0\} \).

\begin{itemize}
    \item \(B\) is called an orthogonal system in \( V \) if all the vectors \( v_i \) are pairwise 
          orthogonal.
    \item An orthogonal system is called an orthonormal system if, in addition, \( \metric{v_i} = 1 \) for 
          all \( i = 1, \ldots, m \).
    \item An orthogonal system that forms a basis of \( V \) is called an orthogonal basis of \( V \).
    \item An orthonormal system that forms a basis of \( V \) is called an orthonormal basis of \( V \).
\end{itemize}

\textbf{Notation:}

\[
    \text{OG-System} \quad \text{OG-Basis} \quad \text{ON-System} \quad \text{ON-Basis}
\]

Using the Kronecker delta symbol \( \delta_{i,j} \):

\[
    \delta_{i,j} =
    \begin{cases}
    1, & \text{if } i = j, \\
    0, & \text{if } i \neq j,
    \end{cases}
\]

we have \( \inner{v_i}{v_j} = \delta_{i,j} \) for any orthonormal system.

\subsection{Writing a vector in terms of the Orthogonal Basis}

\textbf{Example:}

The vectors
\[
    a_1 =
    \begin{bmatrix}
    1 \\
    0 \\
    0
    \end{bmatrix}, \quad
    a_2 =
    \begin{bmatrix}
    0 \\
    1 \\
    0
    \end{bmatrix}
\]

are orthogonal and normalized, but do not form a basis. Thus, they constitute an 
\emph{orthonormal system}.

\textbf{Example:}

The vectors

\[
    a_1 =
    \begin{bmatrix}
    \frac{1}{\sqrt{2}} \\
    \frac{1}{\sqrt{2}} \\
    0
    \end{bmatrix}, \quad
    a_2 =
    \begin{bmatrix}
    -\frac{1}{\sqrt{2}} \\
    \frac{1}{\sqrt{2}} \\
    0
    \end{bmatrix}, \quad
    a_3 =
    \begin{bmatrix}
    0 \\
    0 \\
    1
    \end{bmatrix}
\]
form a basis, are orthogonal and normalized. Therefore, they form an \emph{orthonormal basis}.

\subsection{Linear Independence of Orthogonal Systems}

An orthogonal system \( (v_1, \ldots, v_n) \) is linearly independent.

\textbf{Proof:} 

Let \( \sum_{i=1}^n \lambda_i v_i = 0 \). Then, for any \( v_j \), it follows that

\[
    \inner{\sum_{i=1}^n \lambda_i v_i}{v_j} = \sum_{i=1}^n \lambda_i \inner{v_i}{v_j} = 0.
\]

Due to the orthogonality of the system, all terms vanish except \( \lambda_j \inner{v_j}{v_j} \), 
so

\[
    \lambda_j \inner{v_j}{v_j} = 0.
\]

Since \( v_j \neq 0 \) in any orthogonal system, we have \( \inner{v_j}{v_j} = \metric{v_j}^2 > 0 \), 
implying \( \lambda_j = 0 \) for all \( 1 \leq j \leq n \).

\QED

\subsection{Representation with Respect to an Orthogonal Basis}

Let \( B = (v_1, \ldots, v_n) \) be an orthogonal basis of a vector space \( V \). Then for every 
\( v \in V \), we have:

\[
    v = \sum_{k=1}^n \frac{\inner{v}{v_k}}{\inner{v_k}{v_k}} v_k,
\]

that is, the coordinates of \( v \) with respect to the basis \(B\) are given by

\[
    \left( \frac{\inner{v}{v_k}}{\metric{v_k}^2} \right)_{1 \leq k \leq n}^T.
\]

\subsection{Orthogonal Projection and Direct Sum Decomposition}

Let \( B = (v_1, \ldots, v_m) \) be an orthogonal system in \( V \), and let \( U = L(B) \) be the 
subspace of \( V \) spanned by \(B\).

\begin{itemize}
    \item For every \( v \in V \), the orthogonal projection of \( v \) onto \( U \) is given by:
        
          \[
               p_U(v) = \sum_{i=1}^{m} \frac{\inner{v}{v_i}}{\inner{v_i}{v_i}} v_i.
          \]

    \item Every vector \( v \in V \) can be uniquely written as the sum \( v = p_U(v) + w \), with 
          \( w \in U^\perp \). In this case, we have:
          
          \[
               w = v - p_U(v).
          \]

    \item \( V = U \oplus U^\perp \), which means that every vector in \( V \) can be uniquely 
          decomposed into a sum of a vector from \( U \) and a vector from \( U^\perp \).

    \item If \( \dim(V) = n \), then:
  
          \[
               \dim(U) + \dim(U^\perp) = n \quad \text{for every subspace } U.
          \]

\end{itemize}

\subsection{The Gram-Schmidt Process}

The Gram-Schmidt process is a method for orthonormalizing a set of vectors in an inner product space. 
Given a finite set of linearly independent vectors \( (v_1, \ldots, v_n) \), the process generates an orthonormal basis \( (u_1, \ldots, u_n) \) as follows:

\begin{enumerate}

    \item Set \( u_1 = \frac{v_1}{\metric{v_1}} \).
    
    \item For \( k = 2, \ldots, n \):

    \begin{enumerate}
        
        \item Set \( w_k = v_k - \sum_{j=1}^{k-1} \inner{v_k}{u_j} u_j \).
        
        \item Set \( u_k = \frac{w_k}{\metric{w_k}} \).
    
    \end{enumerate}

    \item The resulting set \( (u_1, \ldots, u_n) \) is an orthonormal basis of the subspace spanned 
          by \( (v_1, \ldots, v_n) \).

        \end{enumerate}

The Gram-Schmidt process can be applied to any finite set of linearly independent vectors in an inner 
product space, and it is particularly useful for constructing orthonormal bases in Euclidean spaces.

\textbf{Proof:}

Let \(W\) be a non-zero finite-dimensional subspace of an inner product space and let 
\((u_1, \dots, u_n)\) be a basis for \(W\). We will show that an \emph{orthonormal basis} exists.

Let \(v_1 = u_1\), then compute the vector \(u_2\) orthogonal to the space \(W_1\) spanned by \(v_1\)

\[
    v_2 = u_2 - proj_{w_1} u_2 = u_2 - \frac{\inner{u_2}{v_1}}{\metric{v_1}} u_2
\]

Then normalize the vector.

We can continue this approach to of projecting the vector of \(u_i\) orthogonal to the space \(W_i\) 
spanned by all basis previously computed \(v_i\) basis vectors and then normalize it 
until we get our orthonormal basis. Note, that because our previous vectors formed a basis we can be 
sure that during this process no zero vector will be generated because then we will have 
some \(u_i\) that is a linear combination of the other vectors. Thus, this would be a contradiction.

\QED

\textbf{Example:}

Let \( v_1 = (1, 0, 0) \), \( v_2 = (1, 1, 0) \), and \( v_3 = (1, 1, 1) \). The Gram-Schmidt process 
yields:

\[
    u_1 = \frac{v_1}{\metric{v_1}} = (1, 0, 0), \quad
    u_2 = \frac{v_2 - \inner{v_2}{u_1} u_1}{\|v_2 - \inner{v_2}{u_1} u_1\|} = 
    \left(0, 1, 0\right), \quad
\]
\[
    u_3 = \frac{v_3 - \inner{v_3}{u_1} u_1 - \inner{v_3}{u_2} u_2}{\|v_3 - \inner{v_3}{u_1} u_1 - \inner{v_3}{u_2} u_2\|} = \left(0, 0, 1\right).
\]

Thus, the orthonormal basis is \( (1, 0, 0), (0, 1, 0), (0, 0, 1) \).

\subsection{Existence of Orthonormal Bases and Orthogonal Projections}

\begin{itemize}

    \item Every finitely generated unitary vector space has an orthonormal basis.
    
    \item This is in general false for vector spaces that are not finitely generated.
    
    \item We now provide the existence proof of the orthogonal projection.
    
    \item Let \( V \) be a unitary vector space and \( U \) a finitely generated subspace. 
          Then for every \( v \in V \), the orthogonal projection \( p_U(v) \) of \( v \) onto \( U \) 
          exists.

    \item Let \( V \) be a finitely generated unitary vector space and \( U \) any subspace. Then we have \( V = U \oplus U^\perp \), and
      
          \[
                \dim(V) = \dim(U) + \dim(U^\perp).
          \]
    
    \item Every hyperplane in \( \R^n \) admits a normal form; the normal vector is 
          unique up to scalar multiplication.
    
    \item Let \( V \) be as above and let \( v_1, \ldots, v_m \in V \). If it is possible to 
          construct orthonormal vectors \( w_1, \ldots, w_m \) from them using the Gram-Schmidt process, 
          then \( (v_1, \ldots, v_m) \) are linearly independent.

\end{itemize}

\subsubsection{Visualization of the Gram-Schmidt process in 2 dimension}

\begin{center}
\begin{tikzpicture}

    \draw[<->] (5,0) -- (-5, 0);
    \draw[<->] (0,5) -- (0, -5);
    \draw[->]  (0,0) -- (2,2)node[below]{\(v1\)};
    \draw[->]  (0,0) -- (1, 3)node[above]{\(v2\)};
    \draw[-, orange] (2,2) -- (1, 3)node[right]{\(v2 - \alpha v1 = w2 \text{ after normalization}\)};
    \draw[->,red]  (0,0) -- (-1, 1)node[below]{\(w2\)};
    \draw[->,blue]  (0,0) -- (1,1)node[below]{\(w1\)};

\end{tikzpicture}
\end{center}

\subsection{Best Approximation}

Let \(V\) be a unitary vector space and \(U\) a finitely generated subspace of \(V\).
A vector \(v* \in U\) is called Best Approximation in \(U\) on \(v\), if:

\[
    \|v* - v\| = \inf_{x \in U}\|x - v\|
\]

also

\[
    \|v* - P_u(v)\| = \inf_{u \in U}\|u - v\|
\]

\textbf{Proof:}

Let \(x \in V\), \(v\) be some point and \(v* = proj_{V} v\). Then 

\[
    v - x = v - v* + v* - x
\]

where \(v - v* \in V^{\bot}\) and \(v* - x \in V\). By the Pythagorean Theorem

\[
    \| x - v \|^2 = \|v* - v\|^2 + \|v* - x\|^2
\]

Therefore

\[
    \|v* - v\|^2 \le \|v* - x\|^2
\]

\QED

